\documentclass[10pt,letterpaper,sans]{moderncv}

\moderncvstyle{banking}
\moderncvcolor{blue}

\usepackage[scale=0.82]{geometry}
\setlength{\hintscolumnwidth}{2.1cm}

\name{Max}{Siegel}
\title{Research Scientist}
\address{Cambridge, MA}{}{}
\email{maxs@mit.edu}

\begin{document}
\makecvtitle

\section{Summary}
\cvitem{}{Research Scientist with deep experience in computational cognitive science, perception, concept learning, and multimodal reasoning. Builds computational models, designs behavioral experiments, and publishes interdisciplinary research spanning human perception, reasoning, and human-AI interaction.}

\section{Experience}
\cventry{2018--Present}{Research Scientist}{MIT, Department of Brain and Cognitive Sciences}{Cambridge, MA}{}{
\begin{itemize}
\item Conduct research in the Computational Cognitive Science group at MIT on perception, concept learning, multimodal reasoning, and human-AI interaction.
\item Develop and evaluate computational accounts of human intelligence using tools from cognitive science, probabilistic modeling, and analysis-by-synthesis.
\item Lead research projects from problem formulation through experiment design, analysis, writing, and publication.
\item Collaborate across psychology, neuroscience, computer science, and developmental science on projects involving vision, sound, language, and social reasoning.
\item Publish work in venues including \textit{Nature Human Behaviour}, \textit{Nature Communications}, \textit{Journal of Experimental Psychology: General}, and the \textit{Proceedings of the Annual Meeting of the Cognitive Science Society}.
\end{itemize}}

\section{Core Strengths}
\cvitem{Research}{Computational cognitive science, visual perception, concept learning, multimodal reasoning, human-AI interaction}
\cvitem{Methods}{Computational modeling, behavioral experiments, statistical analysis, hypothesis-driven research, interdisciplinary collaboration}
\cvitem{Output}{Scientific writing, publication strategy, literature synthesis, presentation of technical work to research audiences}

\section{Selected Publications}
\cvitem{2025}{Machino, Y., \textbf{Siegel, M. H.}, Hofer, M., Tenenbaum, J. B., \& Hawkins, R. Minding the Politeness Gap in Cross-cultural Communication. \textit{Proceedings of the Annual Meeting of the Cognitive Science Society}.}
\cvitem{2023}{Yildirim, I., \textbf{Siegel, M. H.}, Soltani, A. A., Ray Chaudhuri, S., \& Tenenbaum, J. B. Perception of 3D shape integrates intuitive physics and analysis-by-synthesis. \textit{Nature Human Behaviour}.}
\cvitem{2023}{Mazor, M., \textbf{Siegel, M. H.}, \& Tenenbaum, J. B. Prospective search time estimates reveal the strengths and limits of internal models of visual search. \textit{Journal of Experimental Psychology: General}.}
\cvitem{2021}{\textbf{Siegel, M. H.}, Magid, R. W., Pelz, M., Tenenbaum, J. B., \& Schulz, L. E. Children's exploratory play tracks the discriminability of hypotheses. \textit{Nature Communications}.}

\section{Education}
\cventry{2018}{Ph.D. in Brain and Cognitive Sciences}{MIT}{Cambridge, MA}{}{Graduated September 2018.}

\section{Additional Information}
\cvitem{Location}{Cambridge, Massachusetts}
\cvitem{Focus}{Research scientist and applied research roles spanning AI, cognition, perception, and human-centered machine intelligence}

\end{document}
